\documentclass[12pt]{article}
\usepackage[utf8]{inputenc}
\usepackage[T1]{fontenc}
\usepackage{hyperref}
\title{Tama Language Resources for K-12 Curriculum Development}
\author{Compiled by Claude Sonnet 5 (claude-sonnet-5) for Andr\'as Kornai}
\date{}
\begin{document}
\maketitle
\section{Introduction}
Tama (ISO 639-3 \texttt{ten}) is a dormant or extinct Western Tucanoan language of the Caquetá/Orteguaza River region of Colombia, closely associated with Koreguaje. Sources disagree on whether any speakers remain. No curriculum materials exist in search results; only a short historical word list is documented.
\section{Primary Educational Resources}
None found.
\section{Linguistic and Reference Materials}
\begin{itemize}
\item Wikipedia, Tama language (Colombia), including a short word list from an 1882 source: \url{https://en.wikipedia.org/wiki/Tama_language_(Colombia)}
\item Wikipedia, Tucanoan languages (classification): \url{https://en.wikipedia.org/wiki/Tucanoan_languages}
\item Universidad Nacional de Colombia, \emph{Maguar\'e} article discussing Tama's affiliation: \url{https://revistas.unal.edu.co/index.php/maguare/article/download/80-104/15048/42443}
\end{itemize}
\section{Religious and Community Resources}
\begin{itemize}
\item Joshua Project, Tama in Colombia (reports a population of about 600; a mission-oriented source, linguistic claims need caution): \url{https://joshuaproject.net/people_groups/21850/CO}
\end{itemize}
\section{Digital and Audio Resources}
None found.
\section{Curriculum Development Context}
A heritage or revitalization curriculum would have to rely on the closely related Koreguaje language and on community members of Tama descent in Caquetá, after confirming the relationship. Documentation of Tama itself is minimal and should not be over-claimed.
\end{document}